\section{Data and Setting}

The dataset comprises EHR records from pediatric patients (aged 0--18 years) presenting to the SickKids ED between June 2018 and January 2026. After excluding patients who left without being seen ($n = 25{,}397$) and encounters missing the free-text clinical diagnoses required for label imputation ($n = 11{,}396$), the final cohort contained $N = 483{,}705$ unique ED encounters (median age 5.0 years, IQR 2.0--10.0; 54.1\% male).

The model's feature space is limited to information available at triage: demographics (age, sex, weight), vital signs (pulse, respiration, temperature, blood pressure), chief complaint narrative, and Canadian Triage and Acuity Scale (CTAS) score. Post-assessment ICD-10 diagnosis codes, assigned by the treating physician after evaluation, are used only by the label imputation pipeline and excluded from the model's input features to prevent target leakage.

Four binary test-ordered labels define the prediction targets: Arm X-Rays ($n = 1{,}904$ positive encounters), ECGs ($n = 12{,}236$), Appendix Ultrasounds ($n = 3{,}293$), and Testicular Ultrasounds ($n = 2{,}889$). Positive labels are reliable ground truth: they record tests ordered and performed at SickKids. Negative labels are ambiguous. A recorded absence may reflect a genuine clinical decision not to test, or a test performed at an external facility whose records never reach the SickKids EHR. This asymmetric noise structure is the central problem the methods below address.
